% Auto-generated from raw.txt
% Usage:
%   \vocabentry{word}{/ipa/ (pos)}{definition}{example}{note}

\vocabentry{Biodiversity}
{/ˌbaɪoʊdaɪˈvɜːrsəti/ (n)}
{The variety of plant and animal life in the world or in a particular habitat.}
{The loss of biodiversity is one of the most pressing concerns for modern zoologists.}
{Đây là một danh từ dùng để chỉ sự đa dạng sinh học, yếu tố sống còn để duy trì hệ sinh thái.}

\vocabentry{Conservation}
{/ˌkɑːnsərˈveɪʃn/ (n)}
{The protection of plants and animals and natural areas from the damaging effects of human activity.}
{Wildlife conservation groups are working hard to protect the endangered tiger population.}
{Đây là một danh từ dùng để chỉ sự bảo tồn động vật hoang dã và môi trường sống của chúng.}

\vocabentry{Endangered}
{/ɪnˈdeɪndʒərd/ (adj)}
{(Of a species) at serious risk of extinction.}
{The giant panda is perhaps the most famous example of an endangered species.}
{Đây là một tính từ dùng để chỉ các loài đang gặp nguy hiểm, có nguy cơ tuyệt chủng.}

\vocabentry{Habitat}
{/ˈhæbɪtæt/ (n)}
{The natural home or environment of an animal, plant, or other organism.}
{Destruction of the natural habitat is the leading cause of species decline.}
{Đây là một danh từ dùng để chỉ môi trường sống tự nhiên của động vật.}

\vocabentry{Species}
{/ˈspiːʃiːz/ (n)}
{A group of living organisms consisting of similar individuals capable of exchanging genes or interbreeding.}
{Scientists discover hundreds of new insect species in the rainforest every year.}
{Đây là một danh từ dùng để chỉ loài sinh vật.}

\vocabentry{Extinction}
{/ɪkˈstɪŋkʃn/ (n)}
{The state or process of a species, family, or larger group being or becoming extinct.}
{Many prehistoric animals faced extinction due to drastic climate changes.}
{Đây là một danh từ dùng để chỉ sự tuyệt chủng.}

\vocabentry{Zoology}
{/zuˈɑːlədʒi/ (n)}
{The scientific study of the behavior, structure, physiology, classification, and distribution of animals.}
{Advances in zoology have allowed us to understand the complex social structures of elephants.}
{Đây là một danh từ dùng để chỉ ngành động vật học.}

\vocabentry{Ecosystem}
{/ˈiːkoʊsɪstəm/ (n)}
{A biological community of interacting organisms and their physical environment.}
{Every animal plays a specific role in maintaining the balance of the ecosystem.}
{Đây là một danh từ dùng để chỉ hệ sinh thái.}

\vocabentry{Predator}
{/ˈpredətər/ (n)}
{An animal that naturally preys on others.}
{The lion is a top predator in the African savannah.}
{Đây là một danh từ dùng để chỉ động vật ăn thịt hoặc kẻ săn mồi.}

\vocabentry{Prey}
{/preɪ/ (n)}
{An animal that is hunted and killed by another for food.}
{Small rodents are the primary prey for many types of owls.}
{Đây là một danh từ dùng để chỉ con mồi.}

\vocabentry{Mammal}
{/ˈmæml/ (n)}
{A warm-blooded vertebrate animal of a class that is distinguished by the possession of hair or fur and the secretion of milk.}
{Whales and dolphins are marine mammals that spend their entire lives in the water.}
{Đây là một danh từ dùng để chỉ động vật có vú.}

\vocabentry{Vertebrate}
{/ˈvɜːrtɪbrət/ (n)}
{An animal of a large group distinguished by the possession of a backbone or spinal column.}
{Mammals, birds, reptiles, amphibians, and fish are all vertebrates.}
{Đây là một danh từ dùng để chỉ động vật có xương sống.}

\vocabentry{Invertebrate}
{/ɪnˈvɜːrtɪbrət/ (n)}
{An animal lacking a backbone, such as an arthropod, mollusk, or annelid.}
{Insects and spiders make up the majority of the world's invertebrate population.}
{Đây là một danh từ dùng để chỉ động vật không xương sống.}

\vocabentry{Amphibian}
{/æmˈfɪbiən/ (n)}
{A cold-blooded vertebrate animal of a class that comprises the frogs, toads, newts, and salamanders.}
{Amphibians are unique because they can live both on land and in water.}
{Đây là một danh từ dùng để chỉ động vật lưỡng cư.}

\vocabentry{Reptile}
{/ˈreptaɪl/ (n)}
{A vertebrate animal of a class that includes snakes, lizards, crocodiles, turtles, and tortoises.}
{Most reptiles lay eggs and have skin covered in scales or bony plates.}
{Đây là một danh từ dùng để chỉ động vật bò sát.}

\vocabentry{Captivity}
{/kæpˈtɪvəti/ (n)}
{The condition of being imprisoned or confined.}
{Many animal rights activists argue that large mammals should not be kept in captivity.}
{Đây là một danh từ dùng để chỉ tình trạng bị nuôi nhốt, không được sống tự do.}

\vocabentry{Cruelty}
{/ˈkruːəlti/ (n)}
{Callous indifference to or pleasure in causing pain and suffering.}
{Laws against animal cruelty have become much stricter in recent years.}
{Đây là một danh từ dùng để chỉ sự tàn ác, thường dùng trong cụm "animal cruelty" - ngược đãi động vật.}

\vocabentry{Poaching}
{/ˈpoʊtʃɪŋ/ (n)}
{The illegal practice of trespassing on another's authority to hunt or fish.}
{Poaching remains a major threat to the survival of the African elephant.}
{Đây là một danh từ dùng để chỉ nạn săn bắn trộm hoặc đánh bắt trái phép.}

\vocabentry{Domestication}
{/dəˌmestɪˈkeɪʃn/ (n)}
{The process of taming an animal and keeping it as a pet or on a farm.}
{The domestication of dogs occurred thousands of years ago.}
{Đây là một danh từ dùng để chỉ sự thuần hóa động vật.}

\vocabentry{Livestock}
{/ˈlaɪvstɑːk/ (n)}
{Farm animals regarded as an asset.}
{The welfare of livestock is a key issue in the debate over industrial farming.}
{Đây là một danh từ dùng để chỉ gia súc, vật nuôi trong nông nghiệp.}

\vocabentry{Ethology}
{/iˈθɑːlədʒi/ (n)}
{The science of animal behavior.}
{Ethology provides insights into how animals communicate and interact with their environment.}
{Đây là một danh từ dùng để chỉ tập tính học, nghiên cứu hành vi động vật.}

\vocabentry{Migration}
{/maɪˈɡreɪʃn/ (n)}
{Seasonal movement of animals from one region to another.}
{The annual migration of wildebeest is one of nature's greatest spectacles.}
{Đây là một danh từ dùng để chỉ sự di cư của động vật.}

\vocabentry{Hibernation}
{/ˌhaɪbərˈneɪʃn/ (n)}
{The state of inactivity and metabolic depression in endotherms.}
{Bears enter a state of deep hibernation to survive the harsh winter months.}
{Đây là một danh từ dùng để chỉ sự ngủ đông.}

\vocabentry{Camouflage}
{/ˈkæməflɑːʒ/ (n/v)}
{The use of any combination of materials, coloration, or illumination for concealment.}
{The chameleon uses camouflage to blend into its surroundings and hide from predators.}
{Đây là một danh từ/động từ chỉ sự ngụy trang.}

\vocabentry{Sentience}
{/ˈsenʃəns/ (n)}
{The capacity to feel, perceive, or experience subjectively.}
{Many philosophers argue that animal sentience entitles them to moral consideration.}
{Đây là một danh từ chỉ khả năng có cảm giác, tri giác - lập luận chính của quyền động vật.}

\vocabentry{Vivisection}
{/ˌvɪvɪˈsekʃn/ (n)}
{The practice of performing operations on live animals for the purpose of scientific research.}
{Vivisection is a highly controversial practice that faces strong opposition from animal rights groups.}
{Đây là một danh từ chỉ việc giải phẫu động vật sống để nghiên cứu, thường bị chỉ trích.}

\vocabentry{Sanctuary}
{/ˈsæŋktʃueri/ (n)}
{A place where injured or unwanted animals are cared for and protected.}
{The elephant sanctuary provides a home for animals rescued from the circus industry.}
{Đây là một danh từ dùng để chỉ khu bảo tồn hoặc nơi trú ẩn an toàn cho động vật.}

\vocabentry{Endemic}
{/enˈdemɪk/ (adj)}
{(Of a plant or animal) native and restricted to a certain place.}
{The lemur is a primate that is endemic to the island of Madagascar.}
{Đây là một tính từ dùng để chỉ loài đặc hữu, chỉ có ở một khu vực nhất định.}

\vocabentry{Fauna}
{/ˈfɔːnə/ (n)}
{The animals of a particular region, habitat, or geological period.}
{Australia is known for its unique and diverse fauna, including many marsupials.}
{Đây là một danh từ dùng để chỉ quần thể động vật của một vùng.}

\vocabentry{Carnivore}
{/ˈkɑːrnɪvɔːr/ (n)}
{An animal that feeds on flesh.}
{Tigers and wolves are classic examples of carnivores.}
{Đây là một danh từ dùng để chỉ động vật ăn thịt.}

\vocabentry{Herbivore}
{/ˈhɜːrbɪvɔːr/ (n)}
{An animal that feeds on plants.}
{Elephants and giraffes are large herbivores that consume vast amounts of vegetation.}
{Đây là một danh từ dùng để chỉ động vật ăn cỏ.}

\vocabentry{Omnivore}
{/ˈɑːmnɪvɔːr/ (n)}
{An animal or person that eats food of both plant and animal origin.}
{Humans are naturally omnivores, although many choose to follow a vegetarian diet.}
{Đây là một danh từ dùng để chỉ động vật ăn tạp.}

\vocabentry{Nocturnal}
{/nɑːkˈtɜːrnl/ (adj)}
{Done, occurring, or active at night.}
{Owls are nocturnal birds that have excellent vision in low light.}
{Đây là một tính từ dùng để chỉ các loài động vật hoạt động về đêm.}

\vocabentry{Diurnal}
{/daɪˈɜːrnl/ (adj)}
{Active during the day.}
{Most primates are diurnal, meaning they hunt and forage while the sun is up.}
{Đây là một tính từ dùng để chỉ các loài động vật hoạt động vào ban ngày.}

\vocabentry{Evolution}
{/ˌevəˈluːʃn/ (n)}
{The process by which different kinds of living organisms are thought to have developed.}
{Natural selection is the primary mechanism driving biological evolution.}
{Đây là một danh từ dùng để chỉ sự tiến hóa.}

\vocabentry{Adaptation}
{/ˌædæpˈteɪʃn/ (n)}
{The process of change by which an organism or species becomes better suited to its environment.}
{The polar bear's thick layer of fat is an essential adaptation for surviving in the Arctic.}
{Đây là một danh từ dùng để chỉ sự thích nghi.}

\vocabentry{Territory}
{/ˈterətɔːri/ (n)}
{An area of land under the jurisdiction of a ruler or state; an area defended by an animal.}
{Many male birds sing to mark their territory and attract a mate.}
{Đây là một danh từ dùng để chỉ lãnh thổ của động vật.}

\vocabentry{Breeding}
{/ˈbriːdɪŋ/ (n)}
{The production of offspring.}
{The zoo has a successful breeding program for endangered red pandas.}
{Đây là một danh từ dùng để chỉ sự sinh sản hoặc gây giống.}

\vocabentry{Offspring}
{/ˈɔːfsprɪŋ/ (n)}
{A person's child or children; an animal's young.}
{Female crocodiles are known to be very protective of their offspring.}
{Đây là một danh từ dùng để chỉ con cái, thế hệ sau của động vật.}

\vocabentry{Parasite}
{/ˈpærəsaɪt/ (n)}
{An organism that lives in or on another organism (its host) and benefits by deriving nutrients at the host's expense.}
{Ticks and fleas are common parasites that feed on the blood of mammals.}
{Đây là một danh từ dùng để chỉ ký sinh trùng.}

\vocabentry{Host}
{/hoʊst/ (n)}
{An animal or plant on or in which a parasite or commensal organism lives.}
{The parasite can cause serious illness in its host if left untreated.}
{Đây là một danh từ dùng để chỉ vật chủ.}

\vocabentry{Symbiosis}
{/ˌsɪmbaɪˈoʊsɪs/ (n)}
{Interaction between two different organisms living in close physical association, typically to the advantage of both.}
{The relationship between clownfish and sea anemones is a famous example of symbiosis.}
{Đây là một danh từ dùng để chỉ sự cộng sinh.}

\vocabentry{Vertebra}
{/ˈvɜːrtɪbrə/ (n)}
{Each of the series of small bones forming the backbone.}
{The human spinal column is made up of thirty-three vertebrae.}
{Đây là một danh từ dùng để chỉ đốt sống.}

\vocabentry{Physiology}
{/ˌfɪziˈɑːlədʒi/ (n)}
{The branch of biology that deals with the normal functions of living organisms and their parts.}
{Understanding animal physiology is essential for providing proper veterinary care.}
{Đây là một danh từ dùng để chỉ sinh lý học.}

\vocabentry{Morphology}
{/mɔːrˈfɑːlədʒi/ (n)}
{The branch of biology that deals with the form of living organisms, and with relationships between their structures.}
{Scientists study the morphology of fossils to understand how ancient animals looked.}
{Đây là một danh từ dùng để chỉ hình thái học.}

\vocabentry{Genome}
{/ˈdʒiːnoʊm/ (n)}
{The haploid set of chromosomes in a gamete or microorganism.}
{Mapping the honeybee genome has provided clues about their complex social behavior.}
{Đây là một danh từ dùng để chỉ hệ gen.}

\vocabentry{Metabolism}
{/məˈtæbəlɪzəm/ (n)}
{The chemical processes that occur within a living organism in order to maintain life.}
{Reptiles have a much slower metabolism than mammals of the same size.}
{Đây là một danh từ dùng để chỉ sự trao đổi chất.}

\vocabentry{Speciesism}
{/ˈspiːʃizɪzəm/ (n)}
{The assumption of human superiority leading to the exploitation of animals.}
{Animal rights advocates compare speciesism to other forms of prejudice like racism.}
{Đây là một danh từ chỉ sự phân biệt đối xử theo loài, coi con người là thượng đẳng.}

\vocabentry{Welfare}
{/ˈwelfer/ (n)}
{The health, happiness, and fortunes of a person or group; or the well-being of animals.}
{Many consumers now care about animal welfare and choose to buy free-range eggs.}
{Đây là một danh từ chỉ phúc lợi, cụ thể là phúc lợi động vật.}

\vocabentry{Abolitionist}
{/ˌæbəˈlɪʃənɪst/ (n)}
{A person who favors the abolition of a practice or institution.}
{Unlike welfarists, animal rights abolitionists believe that all animal use should stop.}
{Trong quyền động vật, đây là người ủng hộ xóa bỏ hoàn toàn việc sử dụng động vật.}

\vocabentry{Zoologist}
{/zuˈɑːlədʒɪst/ (n)}
{An expert in or student of the behavior, structure, physiology, classification, and distribution of animals.}
{The zoologist spent three years in the field studying the mating habits of penguins.}
{Đây là một danh từ dùng để chỉ nhà động vật học.}

\vocabentry{Taxonomy}
{/tækˈsɑːnəmi/ (n)}
{The branch of science concerned with classification, especially of organisms.}
{Taxonomy allows scientists to organize millions of different species into logical groups.}
{Đây là một danh từ dùng để chỉ phân loại học.}

\vocabentry{Natural selection}
{/ˌnætʃrəl sɪˈlekʃn/ (n)}
{The process whereby organisms better adapted to their environment tend to survive and produce more offspring.}
{Natural selection ensures that the strongest and most adaptable individuals pass on their genes.}
{Đây là một cụm danh từ chỉ sự chọn lọc tự nhiên.}

\vocabentry{Inbreeding}
{/ˈɪnbriːdɪŋ/ (n)}
{Breed from closely related people or animals, especially over many generations.}
{Inbreeding in small populations can lead to an increase in genetic defects.}
{Đây là một danh từ dùng để chỉ sự giao phối cận huyết.}

\vocabentry{Gene pool}
{/ˈdʒiːn puːl/ (n)}
{The stock of different genes in an interbreeding population.}
{A large and diverse gene pool is essential for the long-term survival of a species.}
{Đây là một cụm danh từ dùng để chỉ vốn gen.}

\vocabentry{Hybrid}
{/ˈhaɪbrɪd/ (n)}
{The offspring of two plants or animals of different species or varieties.}
{A mule is a hybrid, the offspring of a male donkey and a female horse.}
{Đây là một danh từ dùng để chỉ cá thể lai.}

\vocabentry{Primate}
{/ˈpraɪmeɪt/ (n)}
{A mammal of an order that includes monkeys, apes, and humans.}
{Chimpanzees are our closest living relatives in the primate family.}
{Đây là một danh từ dùng để chỉ bộ linh trưởng.}

\vocabentry{Rodent}
{/ˈroʊdnt/ (n)}
{A gnawing mammal of an order that includes rats, mice, squirrels, hamsters, and porcupines.}
{Rodents are characterized by a pair of continuously growing incisors in each jaw.}
{Đây là một danh từ dùng để chỉ loài gặm nhấm.}

\vocabentry{Ungulate}
{/ˈʌŋɡjələt/ (n)}
{A hoofed mammal.}
{Horses, cows, and deer are all examples of ungulates.}
{Đây là một danh từ dùng để chỉ động vật có móng guốc.}

\vocabentry{Marsupial}
{/mɑːrˈsuːpiəl/ (n)}
{A mammal of an order whose members are born incompletely developed and are typically carried and suckled in a pouch on the mother's belly.}
{Kangaroos and koalas are the most well-known marsupials from Australia.}
{Đây là một danh từ chỉ thú có túi.}

\vocabentry{Scavenger}
{/ˈskævɪndʒər/ (n)}
{An animal that feeds on carrion, dead plant material, or refuse.}
{Vultures play a vital role in the ecosystem as scavengers, cleaning up animal remains.}
{Đây là một danh từ dùng để chỉ động vật ăn xác thối.}

\vocabentry{Carrion}
{/ˈkæriən/ (n)}
{The decaying flesh of dead animals.}
{Some animals, like hyenas, are both hunters and eaters of carrion.}
{Đây là một danh từ dùng để chỉ xác chết đang phân hủy của động vật.}

\vocabentry{Detritivore}
{/dɪˈtraɪtɪvɔːr/ (n)}
{An animal which feeds on dead organic material, especially plant detritus.}
{Earthworms are detritivores that help to break down leaves and enrich the soil.}
{Đây là một danh từ dùng để chỉ sinh vật ăn vụn hữu cơ.}

\vocabentry{Incubation}
{/ˌɪŋkjuˈbeɪʃn/ (n)}
{The process of incubating eggs, birds, etc.}
{The female bird is responsible for the incubation of the eggs for three weeks.}
{Đây là một danh từ dùng để chỉ sự ấp trứng.}

\vocabentry{Fledgling}
{/ˈfledʒlɪŋ/ (n)}
{A young bird that has just fledged.}
{The fledglings were seen fluttering their wings and trying to leave the nest.}
{Đây là một danh từ dùng để chỉ chim non vừa mới tập bay.}

\vocabentry{Larva}
{/ˈlɑːrvə/ (n)}
{The active immature form of an insect, especially one that differs greatly from the adult and forms the stage between egg and pupa.}
{A caterpillar is the larva stage of a butterfly or moth.}
{Đây là một danh từ dùng để chỉ ấu trùng.}

\vocabentry{Pupa}
{/ˈpjuːpə/ (n)}
{An insect in its inactive immature form between larva and adult.}
{Inside the cocoon, the larva transforms into a pupa before emerging as an adult.}
{Đây là một danh từ dùng để chỉ nhộng.}

\vocabentry{Metamorphosis}
{/ˌmetəˈmɔːrfəsɪs/ (n)}
{The process of transformation from an immature form to an adult form in two or more distinct stages.}
{Frogs undergo a complete metamorphosis from tadpoles to adults.}
{Đây là một danh từ dùng để chỉ sự biến thái - thay đổi hình thái.}

\vocabentry{Pheromone}
{/ˈferəmoʊn/ (n)}
{A chemical substance produced and released into the environment by an animal, especially a mammal or an insect, affecting the behavior or physiology of others of its species.}
{Ants use pheromones to create trails for other members of the colony to follow.}
{Đây là một danh từ chỉ chất dẫn dụ hóa học.}

\vocabentry{Echolocation}
{/ˌekoʊloʊˈkeɪʃn/ (n)}
{The location of objects by reflected sound, in particular that used by animals such as bats and dolphins.}
{Bats use echolocation to navigate and find prey in total darkness.}
{Đây là một danh từ dùng để chỉ sự định vị bằng tiếng vang.}

\vocabentry{Endothermic}
{/ˌendoʊˈθɜːrmɪk/ (adj)}
{(Of an animal) dependent on or capable of the internal generation of heat; warm-blooded.}
{Birds and mammals are endothermic, allowing them to remain active in cold environments.}
{Đây là một tính từ chỉ động vật biến nhiệt/máu nóng.}

\vocabentry{Ectothermic}
{/ˌektoʊˈθɜːrmɪk/ (adj)}
{(Of an animal) dependent on external sources of body heat; cold-blooded.}
{Reptiles are ectothermic and often bask in the sun to raise their body temperature.}
{Đây là một tính từ chỉ động vật đẳng nhiệt/máu lạnh.}

\vocabentry{Invasive species}
{/ɪnˈveɪsɪv ˈspiːʃiːz/ (n)}
{A species that is not native to a specific location and has a tendency to spread to a degree believed to cause damage to the environment.}
{The introduction of invasive species can devastate local wildlife populations.}
{Đây là một cụm danh từ chỉ loài xâm lấn.}

\vocabentry{Feral}
{/ˈferəl/ (adj)}
{(Especially of an animal) in a wild state, especially after escape from captivity or domestication.}
{The island is home to a large population of feral cats that hunt native birds.}
{Đây là một tính từ dùng để chỉ động vật hoang dã, từng được thuần hóa nhưng đã trở lại tự nhiên.}

\vocabentry{Anthropomorphism}
{/ˌænθrəpəˈmɔːrfɪzəm/ (n)}
{The attribution of human characteristics or behavior to a god, animal, or object.}
{Many cartoons use anthropomorphism by giving animals the ability to talk and wear clothes.}
{Đây là một danh từ chỉ nhân hóa động vật.}

\vocabentry{Vivisectionist}
{/ˌvɪvɪˈsekʃənɪst/ (n)}
{A person who performs vivisection.}
{Modern science has moved away from the need for vivisectionists in many areas.}
{Đây là một danh từ dùng để chỉ người thực hiện giải phẫu động vật sống.}

\vocabentry{Bioethics}
{/ˌbaɪoʊˈeθɪks/ (n)}
{The study of ethical issues raised by advances in biology and medicine.}
{The debate over cloning raises many complex bioethical questions.}
{Đây là một danh từ dùng để chỉ đạo đức sinh học.}

\vocabentry{Sentient}
{/ˈsenʃnt/ (adj)}
{Able to perceive or feel things.}
{Animals are sentient beings that deserve to be treated with compassion.}
{Đây là một tính từ dùng để chỉ khả năng có tri giác/cảm xúc.}

\vocabentry{Veganism}
{/ˈviːɡənɪzəm/ (n)}
{The practice of abstaining from the use of animal products, particularly in diet.}
{Veganism is often motivated by a desire to support animal rights.}
{Đây là một danh từ dùng để chỉ chủ nghĩa ăn chay trường/lối sống thuần chay.}

\vocabentry{Factory farming}
{/ˈfæktəri ˈfɑːrmɪŋ/ (n)}
{A system of intensive animal farming developed to maximize profits by using minimal space.}
{Factory farming is often criticized for its poor treatment of animals and environmental impact.}
{Đây là một cụm danh từ chỉ chăn nuôi công nghiệp cường độ cao.}

\vocabentry{Free-range}
{/ˌfriː ˈreɪndʒ/ (adj)}
{(Of livestock or poultry) kept in natural conditions, with freedom of movement.}
{Consumers are increasingly choosing free-range chicken over factory-farmed options.}
{Đây là một tính từ dùng để chỉ phương pháp chăn thả tự nhiên.}

\vocabentry{Endangerment}
{/ɪnˈdeɪndʒərmənt/ (n)}
{The action of putting someone or something at risk or in danger of being harmed or destroyed.}
{Habitat loss is the primary driver of species endangerment worldwide.}
{Đây là một danh từ chỉ tình trạng bị đe dọa tuyệt chủng.}

\vocabentry{Overhunting}
{/ˌoʊvərˈhʌntɪŋ/ (n)}
{Hunting to such an extent that the population of a species is significantly reduced.}
{Overhunting in the 19th century led to the near-extinction of the American bison.}
{Đây là một danh từ dùng để chỉ nạn săn bắn quá mức.}

\vocabentry{Ivory trade}
{/ˈaɪvəri treɪd/ (n)}
{The commercial, often illegal, trade in the tusks of elephants and other animals.}
{The international ban on the ivory trade has helped to slow the decline of elephant populations.}
{Đây là một cụm danh từ chỉ việc buôn bán ngà voi.}

\vocabentry{Trophy hunting}
{/ˈtroʊfi ˈhʌntɪŋ/ (n)}
{The hunting of wild animals for human recreation.}
{Trophy hunting is a highly controversial practice that many people believe should be banned.}
{Đây là một cụm danh từ chỉ săn bắn lấy chiến lợi phẩm.}

\vocabentry{Fur trade}
{/fɜːr treɪd/ (n)}
{The worldwide industry dealing in the acquisition and sale of animal fur.}
{Many high-end fashion brands have committed to stopping their participation in the fur trade.}
{Đây là một cụm danh từ chỉ ngành buôn bán lông thú.}

\vocabentry{Animal testing}
{/ˈænɪml ˈtestɪŋ/ (n)}
{The use of non-human animals in experiments.}
{Many cosmetic companies have replaced animal testing with more humane alternatives.}
{Đây là một cụm danh từ dùng để chỉ thử nghiệm trên động vật.}

\vocabentry{Bycatch}
{/ˈbaɪkætʃ/ (n)}
{The unwanted fish and other marine creatures caught during commercial fishing for a different species.}
{Modern fishing nets are designed to reduce bycatch and protect sea turtles.}
{Đây là một danh từ chỉ sinh vật biển bị đánh bắt nhầm/phụ phẩm đánh bắt.}

\vocabentry{Aquaculture}
{/ˈækwəkʌltʃər/ (n)}
{The rearing of aquatic animals or the cultivation of aquatic plants for food.}
{Aquaculture provides a significant portion of the world's seafood supply.}
{Đây là một danh từ dùng để chỉ ngành nuôi trồng thủy sản.}

\vocabentry{Biodiversity hotspot}
{/ˌbaɪoʊdaɪˈvɜːrsəti ˈhɑːtspɑːt/ (n)}
{A biogeographic region that is both a significant reservoir of biodiversity and is threatened with destruction.}
{The Amazon rainforest is one of the world's most important biodiversity hotspots.}
{Đây là một cụm danh từ chỉ điểm nóng đa dạng sinh học.}

\vocabentry{Gene splicing}
{/ˈdʒiːn ˈsplaɪsɪŋ/ (n)}
{The process by which the DNA of an organism is cut and a gene from another organism is inserted.}
{Gene splicing allows scientists to create genetically modified organisms with specific traits.}
{Đây là một cụm danh từ dùng để chỉ kỹ thuật cắt nối gen.}

\vocabentry{Transgenic}
{/trænzˈdʒenɪk/ (adj)}
{Relating to or denoting an organism that contains genetic material into which DNA from an unrelated organism has been artificially introduced.}
{Transgenic mice are often used in medical research to study human diseases.}
{Đây là một tính từ dùng để chỉ sinh vật biến đổi gen.}

\vocabentry{Cloning}
{/ˈkloʊnɪŋ/ (n)}
{The process of producing genetically identical individuals of an organism.}
{The cloning of Dolly the sheep in 1996 was a major scientific breakthrough.}
{Đây là một danh từ dùng để chỉ sự nhân bản vô tính.}

\vocabentry{Cryopreservation}
{/ˌkraɪoʊprezərˈveɪʃn/ (n)}
{The preservation of cells, tissues, or embryos by cooling them to very low temperatures.}
{Cryopreservation could be used to save the genetic material of endangered species for the future.}
{Đây là một danh từ dùng để chỉ sự bảo quản lạnh/đóng băng sinh học.}

\vocabentry{Biotechnology}
{/ˌbaɪoʊtekˈnɑːlədʒi/ (n)}
{The exploitation of biological processes for industrial and other purposes.}
{Biotechnology has the potential to solve many problems in agriculture and medicine.}
{Đây là một danh từ dùng để chỉ công nghệ sinh học.}

\vocabentry{Ornithology}
{/ˌɔːrnɪˈθɑːlədʒi/ (n)}
{The scientific study of birds.}
{Ornithology has revealed the incredible distances that birds travel during migration.}
{Đây là một danh từ dùng để chỉ điểu học - nghiên cứu về chim.}

\vocabentry{Entomology}
{/ˌentəˈmɑːlədʒi/ (n)}
{The scientific study of insects.}
{Entomology is essential for understanding the role of insects in the ecosystem and in agriculture.}
{Đây là một danh từ dùng để chỉ côn trùng học.}

\vocabentry{Herpetology}
{/ˌhɜːrpəˈtɑːlədʒi/ (n)}
{The branch of zoology concerned with reptiles and amphibians.}
{Herpetology involves studying the unique physiology and behavior of snakes and frogs.}
{Đây là một danh từ dùng để chỉ ngành nghiên cứu bò sát và lưỡng cư.}

\vocabentry{Ichthyology}
{/ˌɪkθiˈɑːlədʒi/ (n)}
{The branch of zoology that deals with fishes.}
{Ichthyology is important for managing fisheries and protecting marine biodiversity.}
{Đây là một danh từ dùng để chỉ ngư học - nghiên cứu về cá.}

\vocabentry{Marine biology}
{/məˈriːn baɪˈɑːlədʒi/ (n)}
{The scientific study of organisms in the ocean or other marine bodies of water.}
{Marine biology exploring the deep sea has led to the discovery of many strange and wonderful creatures.}
{Đây là một cụm danh từ chỉ sinh học biển.}

\vocabentry{Phylogeny}
{/faɪˈlɑːdʒəni/ (n)}
{The branch of biology that deals with phylogenesis.}
{Modern phylogeny uses DNA sequencing to determine the evolutionary relationships between species.}
{Đây là một danh từ dùng để chỉ hệ thống phát sinh loài.}

\vocabentry{Ontogeny}
{/ɑːnˈtɑːdʒəni/ (n)}
{The development of an individual organism or anatomical or behavioral feature from the earliest stage to maturity.}
{Ontogeny focuses on how an animal grows and changes from an embryo to an adult.}
{Đây là một danh từ dùng để chỉ quá trình phát triển cá thể.}

\vocabentry{Naturalist}
{/ˈnætʃrəlɪst/ (n)}
{An expert in or student of natural history.}
{Charles Darwin was a famous naturalist whose work transformed our understanding of the natural world.}
{Đây là một danh từ dùng để chỉ nhà tự nhiên học.}

\vocabentry{Biodiversity net gain}
{/ˌbaɪoʊdaɪˈvɜːrsəti net ɡeɪn/ (n)}
{An approach to development that leaves biodiversity in a better state than before.}
{Developers are increasingly required to demonstrate biodiversity net gain in their projects.}
{Đây là một cụm danh từ chỉ sự gia tăng ròng đa dạng sinh học.}

\vocabentry{Corridor}
{/ˈkɔːrɪdɔːr/ (n)}
{A strip of natural habitat connecting populations of wildlife otherwise separated by human activities.}
{Wildlife corridors allow animals to travel safely between fragmented patches of forest.}
{Đây là một danh từ dùng để chỉ hành lang sinh thái cho động vật.}

\vocabentry{Ecological niche}
{/ˌiːkəˈlɑːdʒɪkl nɪʃ/ (n)}
{The role and position a species has in its environment.}
{No two species can occupy the exact same ecological niche in the same habitat for long.}
{Đây là một cụm danh từ dùng để chỉ ổ sinh thái.}

\vocabentry{Trophic level}
{/ˈtroʊfɪk ˈlevl/ (n)}
{Each of several hierarchical levels in an ecosystem.}
{Energy is lost at each trophic level as it moves up the food chain.}
{Đây là một cụm danh từ dùng để chỉ bậc dinh dưỡng.}

\vocabentry{Food web}
{/ˈfuːd web/ (n)}
{A system of interlocking and interdependent food chains.}
{A complex food web is a sign of a healthy and resilient ecosystem.}
{Đây là một cụm danh từ dùng để chỉ lưới thức ăn.}

\vocabentry{Biomass}
{/ˈbaɪoʊmæs/ (n)}
{The total mass of organisms in a given area or volume.}
{In a healthy ecosystem, the biomass of producers is much larger than that of consumers.}
{Đây là một danh từ dùng để chỉ sinh khối.}

\vocabentry{Sustainable development}
{/səˈsteɪnəbl dɪˈveləpmənt/ (n)}
{Economic development that is conducted without depletion of natural resources.}
{Sustainable development is essential for ensuring the long-term survival of both humans and animals.}
{Đây là một cụm danh từ chỉ phát triển bền vững.}

\vocabentry{Abundance}
{/əˈbʌndəns/ (n)}
{A very large quantity of something.}
{The researchers were surprised by the abundance of rare birds in the area.}
{Trong động vật học, dùng để chỉ sự phong phú về số lượng cá thể trong một quần thể.}

\vocabentry{Aesthetic}
{/esˈθetɪk/ (adj)}
{Concerned with beauty or the appreciation of beauty.}
{Many people support wildlife conservation for both ecological and aesthetic reasons.}
{Sự trân trọng vẻ đẹp tự nhiên của thế giới động vật.}

\vocabentry{Aggregation}
{/ˌæɡrɪˈɡeɪʃn/ (n)}
{A cluster of things that have come or been brought together.}
{The migration of whales involves a massive aggregation of individuals in feeding grounds.}
{Sự tập hợp của các cá thể động vật, ví dụ như đàn cá hoặc đàn chim.}

\vocabentry{Aggression}
{/əˈɡreʃn/ (n)}
{Feelings of anger or antipathy resulting in hostile or violent behavior.}
{Male lions often show aggression when defending their territory from rivals.}
{Hành vi hung hăng của động vật.}

\vocabentry{Altruism}
{/ˈæltruɪzəm/ (n)}
{The belief in or practice of disinterested and selfless concern for the well-being of others.}
{Altruism has been observed in dolphins that help injured members of their pod.}
{Hành vi vị tha ở động vật, ví dụ như con vật cứu giúp đồng loại.}

\vocabentry{Analogous}
{/əˈnæləɡəs/ (adj)}
{(Of structures) performing a similar function but having a different evolutionary origin.}
{The wings of a bird and the wings of a butterfly are analogous structures.}
{Cơ quan tương tự trong giải phẫu so sánh.}

\vocabentry{Anatomy}
{/əˈnætəmi/ (n)}
{The branch of science concerned with the bodily structure of humans, animals, and other living organisms.}
{Students of zoology must have a strong foundation in animal anatomy.}
{Ngành giải phẫu học.}

\vocabentry{Annual}
{/ˈænjuəl/ (adj)}
{Occurring once every year.}
{The annual migration of monarch butterflies is a remarkable natural event.}
{Các sự kiện diễn ra hàng năm như di cư hoặc sinh sản.}

\vocabentry{Anomalous}
{/əˈnɑːmələs/ (adj)}
{Deviating from what is standard, normal, or expected.}
{The bird's anomalous coloration made it easy for predators to spot.}
{Các bất thường về hình thái hoặc hành vi ở động vật.}

\vocabentry{Antagonism}
{/ænˈtæɡənɪzəm/ (n)}
{Active hostility or opposition.}
{There is often strong antagonism between different species of predators in the same area.}
{Sự đối kháng giữa các loài.}

\vocabentry{Apex predator}
{/ˈeɪpeks ˈpredətər/ (n)}
{A predator at the top of a food chain, with no natural predators of its own.}
{Sharks are apex predators that play a crucial role in maintaining marine ecosystems.}
{Kẻ săn mồi đỉnh bảng.}

\vocabentry{Aquarium}
{/əˈkweriəm/ (n)}
{A transparent tank of water in which fish and other water creatures and plants are kept.}
{Public aquariums educate visitors about the importance of ocean conservation.}
{Bể cá hoặc thủy cung.}

\vocabentry{Arboreal}
{/ɑːrˈbɔːriəl/ (adj)}
{(Of animals) living in trees.}
{Monkeys and squirrels are arboreal mammals that have adapted to life in the canopy.}
{Động vật sống trên cây.}

\vocabentry{Archaeology}
{/ˌɑːrkiˈɑːlədʒi/ (n)}
{The study of human history and prehistory through the excavation of sites.}
{Zooarchaeology uses animal remains from archaeological sites to study past human diets.}
{Khảo cổ học giúp tìm hiểu về mối quan hệ cổ xưa giữa người và vật.}

\vocabentry{Articulate}
{/ɑːrˈtɪkjələt/ (adj)}
{Having joints or jointed segments.}
{The articulated limbs of mammals allow for complex movements.}
{Có khớp xương, đặc điểm của động vật có xương sống.}

\vocabentry{Artifact}
{/ˈɑːrtɪfækt/ (n)}
{An object made by a human being.}
{Ancient artifacts include tools made from bone and ivory.}
{Trong quyền động vật, có thể chỉ các công cụ săn bắn cổ.}

\vocabentry{Artificial}
{/ˌɑːrtɪˈfɪʃl/ (adj)}
{Made or produced by human beings rather than occurring naturally.}
{Artificial insemination is often used to maintain genetic diversity in captive populations.}
{Ví dụ: "artificial insemination" - thụ tinh nhân tạo cho động vật quý hiếm.}

\vocabentry{Asexual}
{/eɪˈsekʃuəl/ (adj)}
{(Of reproduction) not involving the fusion of gametes.}
{Some species of sea anemones reproduce through asexual budding.}
{Sinh sản vô tính ở một số loài động vật bậc thấp.}

\vocabentry{Assimilation}
{/əˌsɪməˈleɪʃn/ (n)}
{The process of taking in and fully understanding information or ideas.}
{Digestion is followed by the assimilation of nutrients into the body's tissues.}
{Trong sinh học, đây là sự đồng hóa dưỡng chất.}

\vocabentry{Assumption}
{/əˈsʌmpʃn/ (n)}
{A thing that is accepted as true or as certain to happen, without proof.}
{Zoologists must be careful not to make assumptions about an animal's intelligence based on its appearance.}
{Các giả định trong nghiên cứu hành vi động vật.}

\vocabentry{Atrophy}
{/ˈætrəfi/ (v)}
{(Of body tissue or an organ) waste away, especially as a result of the degeneration of cells.}
{Muscles can atrophy if an animal is confined to a small cage for too long.}
{Sự teo đi của các cơ quan không sử dụng.}

\vocabentry{Attain}
{/əˈteɪn/ (v)}
{Succeed in achieving (something that one has worked for or concentrated on).}
{Some species of turtles can attain a very great age, often over a hundred years.}
{Đạt được kích thước hoặc độ tuổi trưởng thành.}

\vocabentry{Attribute}
{/ˈætrɪbjuːt/ (n)}
{A quality or feature regarded as a characteristic or inherent part of someone or something.}
{Sharp claws and keen eyesight are important attributes for a predator.}
{Thuộc tính hoặc đặc điểm của loài.}

\vocabentry{Auditory}
{/ˈɔːdɪtɔːri/ (adj)}
{Relating to the sense of hearing.}
{Dogs have an incredibly acute auditory sense, allowing them to hear high-pitched sounds.}
{Thuộc về thính giác.}

\vocabentry{Autonomy}
{/ɔːˈtɑːnəmi/ (n)}
{The right or condition of self-government.}
{Some animal rights theorists believe that sentient animals possess a form of moral autonomy.}
{Quyền tự chủ của động vật trong triết học quyền động vật.}

\vocabentry{Avian}
{/ˈeɪviən/ (adj)}
{Relating to birds.}
{Avian influenza is a highly contagious viral disease that affects many bird species.}
{Thuộc về các loài chim.}

\vocabentry{Avoidance}
{/əˈvɔɪdəns/ (n)}
{The action of keeping away from or not doing something.}
{Predator avoidance is a key survival strategy for many small animals.}
{Hành vi tránh né kẻ thù.}

\vocabentry{Awareness}
{/əˈwernəs/ (n)}
{Knowledge or perception of a situation or fact.}
{There is a growing awareness of the suffering caused by the illegal wildlife trade.}
{Sự nhận thức của công chúng về quyền động vật.}

\vocabentry{Backbone}
{/ˈbækboʊn/ (n)}
{The series of vertebrae extending from the skull to the pelvis, forming the main support of the body.}
{The backbone protects the spinal cord and provides structural support for the body.}
{Xương sống.}

\vocabentry{Balance}
{/ˈbæləns/ (n)}
{A condition in which different elements are equal or in the correct proportions.}
{We must work to restore the natural balance that has been disturbed by human activity.}
{Sự cân bằng sinh thái.}

\vocabentry{Banish}
{/ˈbænɪʃ/ (v)}
{Send (someone) away from a country or place as an official punishment.}
{In some animal societies, weak or aggressive individuals are banished from the group.}
{Đuổi một cá thể ra khỏi đàn.}

\vocabentry{Barrier}
{/ˈbæriər/ (n)}
{A fence or other obstacle that prevents movement or access.}
{Roads and fences can act as barriers that prevent animals from migrating safely.}
{Rào cản ngăn cách môi trường sống.}

\vocabentry{Bask}
{/bæsk/ (v)}
{Lie exposed to warmth and light, typically from the sun, for relaxation or pleasure.}
{Lizards spend much of the morning basking on rocks to warm their bodies.}
{Sưởi nắng, đặc biệt ở bò sát.}

\vocabentry{Beak}
{/biːk/ (n)}
{A bird's horny projecting jaws; a bill.}
{Different species of birds have beaks that are adapted to their specific diet.}
{Mỏ chim.}

\vocabentry{Bear}
{/ber/ (v)}
{(Of a female animal) give birth to (young).}
{Female mammals bear live young and provide them with milk.}
{Sinh con.}

\vocabentry{Beast}
{/biːst/ (n)}
{An animal, especially a large or dangerous four-footed one.}
{The ancient myths tell stories of brave heroes battling ferocious beasts.}
{Thú vật, thường là loài lớn hoặc nguy hiểm.}

\vocabentry{Behavior}
{/bɪˈheɪvjər/ (n)}
{The way in which an animal or person acts in response to a particular situation or stimulus.}
{Zoologists spend many hours observing the social behavior of primates.}
{Hành vi động vật.}

\vocabentry{Beneficiary}
{/ˌbenɪˈfɪʃieri/ (n)}
{A person who derives advantage from something.}
{Endangered species are the primary beneficiaries of wildlife protection laws.}
{Động vật là đối tượng hưởng lợi từ các chính sách bảo tồn.}

\vocabentry{Bilateral}
{/ˌbaɪˈlætərəl/ (adj)}
{Having or relating to two sides.}
{Most animals exhibit bilateral symmetry, meaning their left and right sides are mirror images.}
{Đối xứng hai bên, một đặc điểm cơ thể của nhiều loài động vật.}

\vocabentry{Binary}
{/ˈbaɪnəri/ (adj)}
{Relating to, composed of, or involving two things.}
{Bacteria reproduce through binary fission, where one cell divides into two identical ones.}
{Trong di truyền học, ví dụ: binary fission - phân đôi.}

\vocabentry{Bioaccumulation}
{/ˌbaɪoʊəˌkjuːmjuˈleɪʃn/ (n)}
{The accumulation of a substance, such as a toxic chemical, in various tissues of a living organism.}
{Bioaccumulation of mercury in fish can pose a serious threat to human health.}
{Sự tích tụ sinh học của chất độc.}

\vocabentry{Biogeography}
{/ˌbaɪoʊdʒiˈɑːɡrəfi/ (n)}
{The branch of biology that deals with the geographical distribution of plants and animals.}
{Biogeography helps us understand why certain species are found only in specific regions.}
{Địa sinh học.}

\vocabentry{Biological}
{/ˌbaɪəˈlɑːdʒɪkl/ (adj)}
{Relating to biology or living organisms.}
{The study of biological diversity is essential for protecting the planet's ecosystems.}
{Thuộc về sinh học.}

\vocabentry{Biome}
{/ˈbaɪoʊm/ (n)}
{A large naturally occurring community of flora and fauna occupying a major habitat.}
{The tropical rainforest is a biome characterized by high rainfall and warm temperatures.}
{Khu sinh học.}

\vocabentry{Bipedal}
{/ˌbaɪˈpiːdl/ (adj)}
{(Of an animal) using only two legs for walking.}
{Humans and birds are bipedal, while most other mammals are quadrupedal.}
{Di chuyển bằng hai chân.}

\vocabentry{Birth}
{/bɜːrθ/ (n)}
{The emergence of a young animal or human from the body of its mother.}
{The zoo celebrated the birth of a rare white rhino calf last week.}
{Sự sinh đẻ.}

\vocabentry{Blubber}
{/ˈblʌbər/ (n)}
{The fat of sea mammals, especially whales and seals.}
{Blubber provides insulation and energy storage for animals living in cold oceans.}
{Lớp mỡ dưới da của động vật biển.}

\vocabentry{Bond}
{/bɑːnd/ (n/v)}
{A strong force of attraction; to join or be joined.}
{The bond between a mother and her young is incredibly strong in many species.}
{Sự gắn kết giữa các con vật trong đàn.}

\vocabentry{Botany}
{/ˈbɑːtəni/ (n)}
{The scientific study of plants.}
{Zoology and botany are two major branches of the biological sciences.}
{Nghiên cứu thực vật thường đi đôi với động vật học.}

\vocabentry{Bounty}
{/ˈbaʊnti/ (n)}
{A sum paid for killing or capturing a person or animal.}
{In the past, governments offered bounties for the pelts of wolves and cougars.}
{Tiền thưởng cho việc săn bắt động vật gây hại - nay bị chỉ trích.}

\vocabentry{Brackish}
{/ˈbrækɪʃ/ (adj)}
{(Of water) slightly salty, as is the mixture of river water and seawater in estuaries.}
{Many species of fish are adapted to living in the brackish water of mangrove swamps.}
{Nước lợ, môi trường sống của một số loài cá.}

\vocabentry{Branch}
{/bræntʃ/ (n)}
{A part of a tree which grows out from the trunk.}
{Mammals represent one major branch of the vertebrate family tree.}
{Trong phân loại, là một nhánh trong cây tiến hóa.}

\vocabentry{Breed}
{/briːd/ (n/v)}
{A stock of animals within a species having a distinctive appearance; to produce offspring.}
{The German Shepherd is a popular breed of dog known for its intelligence and loyalty.}
{Giống vật nuôi hoặc hoạt động sinh sản.}

\vocabentry{Broad}
{/brɔːd/ (adj)}
{Covering a large area; general.}
{The species has a broad distribution across most of North America.}
{Dải phân bố rộng của một loài.}

\vocabentry{Brood}
{/bruːd/ (n/v)}
{A family of young animals, especially birds; to sit on eggs to hatch them.}
{The mother hen spent all day caring for her large brood of chicks.}
{Lứa con hoặc việc ấp trứng.}

\vocabentry{Burrow}
{/ˈbɜːroʊ/ (n/v)}
{A hole or tunnel dug by a small animal; to make a hole or tunnel.}
{Rabbits live in complex underground burrows called warrens.}
{Hang động vật hoặc việc đào hang.}

\vocabentry{Calamity}
{/kəˈlæməti/ (n)}
{An event causing great and often sudden damage or distress; a disaster.}
{A forest fire is a major calamity that can destroy entire animal populations.}
{Thảm họa thiên nhiên ảnh hưởng đến động vật.}

\vocabentry{Calf}
{/kæf/ (n)}
{The young of certain mammals, such as cows, elephants, and whales.}
{The baby elephant, or calf, follows its mother closely for the first few years of its life.}
{Con non của một số loài thú lớn.}

\vocabentry{Call}
{/kɔːl/ (n/v)}
{A cry made by an animal; to make such a cry.}
{The distinct call of the loon is a familiar sound in many northern lakes.}
{Tiếng gọi hoặc tiếng kêu của động vật.}

\vocabentry{Canine}
{/ˈkeɪnaɪn/ (adj/n)}
{Relating to or resembling a dog or dogs.}
{Wolves and foxes are members of the canine family.}
{Thuộc về họ chó hoặc răng nanh.}

\vocabentry{Canopy}
{/ˈkænəpi/ (n)}
{The uppermost branches of the trees in a forest.}
{Many arboreal animals spend their entire lives in the high rainforest canopy.}
{Tầng tán rừng, nơi cư trú của nhiều loài động vật.}

\vocabentry{Capability}
{/ˌkeɪpəˈbɪləti/ (n)}
{Power or ability.}
{Some species of birds have the capability to mimic human speech.}
{Khả năng của động vật.}

\vocabentry{Capacity}
{/kəˈpæsəti/ (n)}
{The maximum amount that something can contain.}
{The national park has reached its carrying capacity for the local deer population.}
{Sức chứa của môi trường sống.}

\vocabentry{Carcass}
{/ˈkɑːrkəs/ (n)}
{The dead body of an animal.}
{Scavengers like vultures feed on the carcasses of animals that have died of natural causes.}
{Xác động vật.}

\vocabentry{Carnivorous}
{/kɑːrˈnɪvərəs/ (adj)}
{(Of an animal) feeding on other animals.}
{The Venus flytrap is an unusual carnivorous plant that traps and eats insects.}
{Có tính chất ăn thịt.}

\vocabentry{Category}
{/ˈkætəɡɔːri/ (n)}
{A class or division of people or things regarded as having particular shared characteristics.}
{Animals can be organized into different categories based on their diet or habitat.}
{Hạng mục phân loại động vật.}

\vocabentry{Cattle}
{/ˈkætl/ (n)}
{Large ruminant animals with horns and cloven hoofs, domesticated for meat or milk.}
{Cattle farming is a major industry in many parts of the world.}
{Gia súc lớn, trâu bò.}

\vocabentry{Cavern}
{/ˈkævərn/ (n)}
{A cave, or a chamber in a cave, typically a large one.}
{The deep cavern provided a cool and dark shelter for the colony of bats.}
{Hang động lớn, nơi trú ẩn của dơi hoặc gấu.}

\vocabentry{Cell}
{/sel/ (n)}
{The smallest structural and functional unit of an organism.}
{Every living animal is made up of millions of specialized cells.}
{Tế bào.}

\vocabentry{Census}
{/ˈsensəs/ (n)}
{An official count or survey of a population.}
{Zoologists conduct a regular census of the tiger population to monitor their numbers.}
{Điều tra số lượng cá thể trong quần thể động vật.}

\vocabentry{Characteristic}
{/ˌkærəktəˈrɪstɪk/ (n/adj)}
{A feature or quality belonging typically to a person, place, or thing.}
{The long neck is a well-known characteristic of the giraffe.}
{Đặc điểm đặc trưng của loài.}

\vocabentry{Charge}
{/tʃɑːrdʒ/ (n/v)}
{A headlong rush by an animal or group.}
{The rhino made a sudden charge at the safari jeep.}
{Sự tấn công, lao tới của động vật.}

\vocabentry{Chitin}
{/ˈkaɪtɪn/ (n)}
{A fibrous substance consisting of polysaccharides, forming the major constituent in the exoskeleton of arthropods.}
{The hard outer shell of an insect is made of a material called chitin.}
{Chất kitin, tạo nên vỏ của động vật chân đốt.}

\vocabentry{Chordate}
{/ˈkɔːrdeɪt/ (n)}
{An animal of the large phylum Chordata, comprising the vertebrates together with the sea squirts and lancelets.}
{Chordates are characterized by having a notochord at some stage of their development.}
{Động vật có dây sống.}

\vocabentry{Chronic}
{/ˈkrɑːnɪk/ (adj)}
{(Of an illness) persisting for a long time or constantly recurring.}
{Older dogs often suffer from chronic joint pain and arthritis.}
{Đây là một tính từ dùng để chỉ các bệnh mãn tính, kéo dài và khó chữa dứt điểm.}

\vocabentry{Circulation}
{/ˌsɜːrkjəˈleɪʃn/ (n)}
{Movement to and fro or around something, especially that of fluid in a closed system.}
{The heart is responsible for the circulation of blood throughout the animal's body.}
{Sự tuần hoàn máu ở động vật.}

\vocabentry{Claw}
{/klɔː/ (n/v)}
{A curved pointed horny nail on each digit of the foot in birds, lizards, and some mammals.}
{Eagles use their sharp claws to grab and hold onto their prey.}
{Móng vuốt hoặc hành động cào.}

\vocabentry{Cleft}
{/kleft/ (adj)}
{Split, divided, or partially divided into two.}
{Cows and pigs have cleft hoofs, which are divided into two toes.}
{Ví dụ: "cleft hoof" - móng chẻ ở gia súc.}

\vocabentry{Climate}
{/ˈklaɪmət/ (n)}
{The weather conditions prevailing in an area in general or over a long period.}
{Many species are struggling to adapt to the rapidly changing global climate.}
{Khí hậu ảnh hưởng đến sự phân bố động vật.}

\vocabentry{Climax}
{/ˈklaɪmæks/ (n)}
{The most intense, exciting, or important point of something.}
{An old-growth forest is an example of a climax community in a stable ecosystem.}
{Trong sinh thái là "climax community" - quần xã đỉnh cực.}

\vocabentry{Coexistence}
{/ˌkoʊɪɡˈzɪstəns/ (n)}
{The state or condition of living together at the same time or in the same place.}
{We need to find ways to promote the peaceful coexistence of humans and large predators.}
{Sự chung sống hòa bình giữa người và động vật.}

\vocabentry{Cold-blooded}
{/ˌkoʊld ˈblʌdɪd/ (adj)}
{(Of an animal) having a body temperature that varies with that of the environment.}
{Snakes are cold-blooded animals and need to sunbathe to warm up.}
{Động vật máu lạnh.}

\vocabentry{Colony}
{/ˈkɑːləni/ (n)}
{A group of animals of one kind living together.}
{Thousands of seabirds live in a large colony on the cliff face.}
{Thuộc địa hoặc bầy đàn tập trung.}

\vocabentry{Coloration}
{/ˌkʌləˈreɪʃn/ (n)}
{The appearance of something, especially an animal or plant, in terms of its colors.}
{The bright coloration of some frogs serves as a warning to predators.}
{Màu sắc cơ thể.}

\vocabentry{Commensalism}
{/kəˈmensəlɪzəm/ (n)}
{An association between two organisms in which one benefits and the other derives neither benefit nor harm.}
{Remora fish attached to sharks are a classic example of commensalism.}
{Hội sinh.}

\vocabentry{Community}
{/kəˈmjuːnəti/ (n)}
{A group of people or animals living in the same place or having a particular characteristic in common.}
{A biological community includes all the different species living in a particular area.}
{Quần xã sinh vật.}

\vocabentry{Compassion}
{/kəˈpæʃn/ (n)}
{Sympathetic pity and concern for the sufferings or misfortunes of others.}
{Animal rights activists believe that we should treat all living beings with compassion.}
{Lòng trắc ẩn đối với động vật.}

\vocabentry{Competition}
{/ˌkɑːmpəˈtɪʃn/ (n)}
{Interaction between organisms in which both are harmed.}
{There is intense competition for food and water in the African savannah during a drought.}
{Sự cạnh tranh nguồn sống giữa các loài.}

\vocabentry{Complex}
{/ˈkɑːmpleks/ (adj)}
{Consisting of many different and connected parts.}
{The social structures of honeybees are incredibly complex and highly organized.}
{Tính phức tạp của hệ thống sinh học hoặc hành vi.}

